% !PRINTED 344
\weberpage{344}

Endlich bilden wir noch je ein $(0,\,2)$ und $(2,\,0)$ enthaltendes
Tripel
\[
(0,\,2)\ (1,\,2)\ (2,\,2), \quad (2,\,0)\ (2,\,1)\ (2,\,2),
\]
womit alle zwölf Tripel, und damit die Anordnung \S.\,91, (20)
bestimmt sind.

\S.\,93.

\gesperrt{Die Gruppe der Tripelgleichungen.}

Es soll nun die \gesperrt{Galois}'sche Gruppe $P$ der Tripelgleichungen
bestimmt werden. Zunächst ist klar, dass in dieser Gruppe nur
solche Permutationen vorkommen, bei denen jedes Tripel wieder
in ein Tripel übergeht. Denn nehmen wir an, die Wurzeln eines
Tripels $x_1, x_2, x_3$ gehen durch eine Permutation in $y_1, y_2, y_3$
über, so folgt, da diese Permutation auf die Gleichung (1) (\S.\,92)
anwendbar ist,
\[
y_1 = \Theta\,(y_2, y_3),
\]
und folglich ist $y_1$ die dritte Wurzel des durch $y_2, y_3$ bestimmten
Tripels.

Gehen wir nun zu der Bezeichnung für die Wurzeln:
\weq{1}{x = (\xi, \eta)}
über, so ergiebt sich aus der Zusammenstellung der Tripel in
den §§.\,91 und 92, dass drei Wurzeln
\weq{2}{x_1 = (\xi_1, \eta_1), \quad x_2 = (\xi_2, \eta_2), \quad x_3 = (\xi_3, \eta_3)}
immer dann und nur dann ein Tripel bilden, wenn
\weq{3}{\begin{array}{l}
\xi_1 + \xi_2 + \xi_3 \equiv 0\\
\eta_1 + \eta_2 + \eta_3 \equiv 0
\end{array} \quad (\mathrm{mod}\ 3).}
Nennen wir eine Permutationsgruppe der $(\xi, \eta)$, deren Substitutionen
 alle der vollständigen linearen Congruenzgruppe für
den Modul 3 angehören (\S.\,78), eine \gesperrt{lineare Gruppe}, so gilt
der Satz:

\gesperrt{1. Die Gruppe $P$ einer Tripelgleichung ist immer
linear.}

Nach \S.\,77 lässt sich jede Permutation der Grössen $(\xi, \eta)$
überhaupt in der Form darstellen:
\weq{4}{\xi' \equiv \varphi\,(\xi, \eta), \quad \eta' \equiv \psi\,(\xi, \eta) \quad (\mathrm{mod}\ 3),}
%CONT>>
